% !PRINTED 405
\weberpage{405}
the decision of our question assume $\nu\leqq\frac{1}{2}n$. Moreover it
follows from \S.\,2 that the coefficients $\alpha$ and $\beta$ of $\varphi$ and $\psi$
must be \gesperrt{integers}.

Now in order to decide on the possibility of such a division, one
takes the function $\varphi(x)$ with indeterminate
coefficients $\alpha$, and divides $f(x)$ by $\varphi(x)$ (according to \S.\,3). There
then results a remainder of degree $\nu-1$, and if one sets its
coefficients, which are all rational functions of the $\alpha$,
equal to zero, one obtains $\nu$ equations which these
coefficients must satisfy. By elimination (\S.\,50) one can
produce an equation which contains only one of these unknowns,
and which must then have an integral root. It
will often be more expedient, however, not actually to carry out this
elimination,. but to attempt directly to satisfy the system of the % [sic] stray comma-period as printed
$\nu$ equations by integral values of the $\alpha$. The
selection of the numbers to be tried can be greatly restricted by many
an artifice resting on theorems of number theory.

As an example we shall take an equation of the tenth degree
from the theory of elliptic functions, where
such problems very frequently occur in such a way that the
decomposability into factors of a definite degree is theoretically established,
and where it is then a matter of finding these
factors $^1$). Such an equation is
\[
x^{10} - 36\,x^8 + 528\,x^6 - 3897\,x^4 + 14354\,x^2 - 21025 = 0.
\]
We ask whether the left side is decomposable into two factors of the
fifth degree. If one of the two factors is
\[
x^5 + \alpha x^4 + \beta x^3 + \gamma x^2 + \delta x + \varepsilon
\]
then, since no odd powers occur in the given equation,
the other is
\[
x^5 - \alpha x^4 + \beta x^3 - \gamma x^2 + \delta x - \varepsilon,
\]
and hence we must have:
\[
\begin{aligned}
x^{10} - 36\,x^8 + 528\,x^6 - 3897\,x^4 + 14354\,x^2 - 21025 \\
= (x^5 + \beta x^3 + \delta x)^2 - (\alpha x^4 + \gamma x^2 + \varepsilon)^2.
\end{aligned}
\]

\noindent\rule{4cm}{0.4pt}

$^1$) Cf.\ the author's work: ``Elliptische Functionen und algebraische
Zahlen'', Braunschweig 1891. The example treated here, from the theory
of the transformation of degree 47, is taken from the paper: ``Ein
Beitrag zur Transformationstheorie der elliptischen Functionen etc.'' by
\gesperrt{H.\ Weber} (1893). Mathematische Annalen, Bd.\ 43.
